\chapter{Context Assembly (Token Budgeting \& Speculative Decoding)}

Cursor is widely considered the best AI IDE in the world. Its secret is the \textbf{Orchestrator Backend}, which dynamically assembles a context window optimized to the exact token limit.

\section{The Cursor Backend Architecture}

Below is the scaled diagram of the context routing loop, showing how inputs flow through the prioritizer to feed the model.

\begin{figure}[h]
\centering
\begin{tikzpicture}[
    scale=0.8, transform shape,
    node distance=0.6cm and 1.2cm,
    box/.style={draw, rectangle, rounded corners, minimum height=0.7cm, align=center, fill=blue!10},
    arrow/.style={thick,->,>=stealth}
]

% Client Group
\node (file) [box] {Active File Editor};
\node (linter) [box, below=of file] {Diagnostics/Linter};
\node (term) [box, below=of linter] {Terminal Output};

\begin{pgfonlayer}{background}
\node[draw, dashed, inner sep=0.2cm, fill=gray!3, fit=(file) (linter) (term), label=above:\textbf{Client (IDE)}] (client) {};
\end{pgfonlayer}

% Orchestrator Group
\node (budget) [box, fill=yellow!20, right=of linter, xshift=1.5cm] {Token Budget Allocator};
\node (rag) [box, fill=green!20, above=of budget] {RAG / Fuzzy Search};
\node (compiler) [box, fill=purple!20, below=of budget] {Prompt Compiler};

\begin{pgfonlayer}{background}
\node[draw, dashed, inner sep=0.2cm, fill=gray!3, fit=(rag) (budget) (compiler), label=above:\textbf{Orchestrator}] (orch) {};
\end{pgfonlayer}

% LLM Group
\node (llm) [box, fill=red!20, right=of budget, xshift=1.5cm] {Claude 3.5 / GPT-4o};

% Arrows
\draw[arrow] (file.east) -- (budget.west);
\draw[arrow] (linter.east) -- (budget.west);
\draw[arrow] (term.east) -- (budget.west);
\draw[arrow] (rag.south) -- (budget.north);
\draw[arrow] (budget.south) -- node[right, font=\footnotesize] {Prune} (compiler.north);
\draw[arrow] (compiler.east) -- node[above, font=\footnotesize] {Prompt} (llm.west);

\end{tikzpicture}
\caption{Cursor Backend Context Assembly}
\end{figure}

\section{Fuzzy Search vs. Embeddings in IDEs}

RAG systems in IDEs cannot rely solely on semantic embeddings.
\begin{itemize}
    \item \textbf{Embeddings:} Excel at conceptual search.
    \item \textbf{Fuzzy Keyword Search (BM25 / Levenshtein):} Excel at exact-name searches (e.g. searching ``updateUserTokenV3'').
\end{itemize}

Production IDEs use a hybrid scoring system. They compute the semantic score ($S_{\text{semantic}}$) and a fuzzy keyword match score ($S_{\text{keyword}}$), then blend them:
\begin{equation}
S_{\text{final}} = \alpha S_{\text{semantic}} + (1 - \alpha) S_{\text{keyword}}
\end{equation}
where $\alpha \approx 0.6$. The Levenshtein distance $D(A,B)$ measures the edit distance between strings and is calculated using a dynamic programming matrix:
\begin{equation}
D(i,j) = \min \begin{cases}
  D(i-1, j) + 1 \\
  D(i, j-1) + 1 \\
  D(i-1, j-1) + \text{cost}
\end{cases}
\end{equation}

\section{Sliding Window Context \& Truncation Algorithms}

When agent loops make successive roundtrips, context accumulates. If the token count exceeds the engine limits, older history must be pruned. Runtimes execute a **Sliding Window eviction queue**, prioritizing system files and diagnostics.


\begin{center}
\noindent\rule{0.6\textwidth}{0.4pt} \\
\vspace{0.3em}
\textit{[Code implementation is available in the companion coding handbook: \texttt{coding-handbook/ch07\_context\_assembly/}]} \\
\noindent\rule{0.6\textwidth}{0.4pt}
\end{center}


\section{Speculative Decoding (Fast Diffs)}

Because 90\% of a code edit involves rewriting identical lines of code, a smaller ``draft'' model guesses the next 10 tokens instantly. The large model then verifies all 10 tokens in a single forward pass. If they match, 10 tokens are generated simultaneously. This decreases code patch compilation latency significantly.

\section{MemGPT-Style Episodic Memory Paging}

While simple sliding windows prevent out-of-token errors, they destroy historical relationships. Production agents require long-term state awareness, which is achieved by modeling memory as an Operating System memory hierarchy.

\subsection{Memory Tiers}
The agent's memory is divided into three logical tiers:
\begin{itemize}
    \item \textbf{Core Memory (RAM):} Placed directly inside the system prompt window. Contains user schemas and the agent's current state. Hard-budgeted (e.g., 2,000 tokens).
    \item \textbf{Recall Memory (Swap Space):} Vector-indexed logs of all past chats. Queried on-demand.
    \item \textbf{Archival Memory (Disk):} Cold, key-value flat files containing static rules and domain facts.
\end{itemize}

\subsection{Active Memory Operations}
Instead of relying on the system to retrieve facts automatically, the agent acts as the memory controller. It executes specialized tool functions to modify its own RAM:
\begin{align}
\text{edit\_core\_memory(key, val)} &\implies \text{Updates value in system prompt window.} \\
\text{search\_recall\_memory(query)} &\implies \text{Searches database chat logs.} \\
\text{write\_to\_archival(key, fact)} &\implies \text{Saves long-term cold facts.}
\end{align}

\begin{center}
\noindent\rule{0.6\textwidth}{0.4pt} \\
\vspace{0.3em}
\textit{[Code implementation is available in the companion coding handbook: \texttt{coding-handbook/ch07\_context\_assembly/episodic\_paged\_memory.py}]} \\
\noindent\rule{0.6\textwidth}{0.4pt}
\end{center}

